% LuaLaTeX document — compile with: lualatex inter_voxel_propagator.tex
\documentclass[11pt,a4paper]{article}

\usepackage{fontspec}
\setmainfont{Latin Modern Roman}
\setmonofont{Latin Modern Mono}

\usepackage{amsmath,amssymb,amsthm,bm}
\usepackage{booktabs}
\usepackage{geometry}
\geometry{margin=2.5cm}
\usepackage{hyperref}
\hypersetup{colorlinks=true,linkcolor=blue!60!black,citecolor=green!50!black,urlcolor=blue!70!black}
\usepackage{enumitem}
\usepackage{xcolor}
\usepackage{mathtools}
\usepackage{array}
\usepackage{longtable}

\newcommand{\bR}{\bm{R}}
\newcommand{\bk}{\bm{k}}
\newcommand{\bw}{\bm{w}}
\newcommand{\bs}{\bm{s}}
\newcommand{\kh}{\hat{\bm{k}}}
\newcommand{\order}{\mathcal{O}}
\newcommand{\D}{\Delta}
\newcommand{\tr}{\operatorname{tr}}
\newcommand{\diag}{\operatorname{diag}}

\theoremstyle{plain}
\newtheorem{theorem}{Theorem}
\newtheorem{proposition}[theorem]{Proposition}
\theoremstyle{remark}
\newtheorem{remark}[theorem]{Remark}

\title{Analytical Inter-Voxel Strain Propagator\\
\large Complete Derivation and Numerical Tables}
\author{Research Notes}
\date{April 2026}

\begin{document}
\maketitle
\tableofcontents
\vspace{1em}

%======================================================================
\section{Introduction}
\label{sec:intro}
%======================================================================

The Foldy--Lax multiple-scattering equation
\begin{equation}
  \mathbf{T} = \mathbf{T}_0(\mathbf{I} - \mathbf{G}_0\mathbf{T}_0)^{-1}
  \label{eq:foldy-lax}
\end{equation}
relates the full T-matrix~$\mathbf{T}$ of a collection of scatterers
to the single-site T-matrix~$\mathbf{T}_0$ and the inter-site
propagator~$\mathbf{G}_0$.  For cubic voxels on a lattice, the
propagator couples volume-averaged displacement and strain degrees of
freedom between sites separated by lattice vector~$\bR$.

The inter-site propagator~$P_{ijkl}(\bR,\omega)$ is the
volume-averaged Green's tensor second derivative:
\begin{equation}
  P_{ijkl}(\bR,\omega) \;=\;
  \int_A\!\int_B
  \frac{\partial^2 G_{ik}(\bs-\bs'+\bR,\,\omega)}%
       {\partial x_j\,\partial x_l}\,d^3s\,d^3s'.
  \label{eq:Pijkl-def}
\end{equation}
In Fourier space:
\begin{equation}
  P_{ijkl}(\bR,\omega) \;=\;
  \int\!\frac{d^3k}{(2\pi)^3}\;
  \bigl|\tilde{f}_0(\bk)\bigr|^2\,
  \tilde{\Gamma}_{ijkl}(\bk,\omega)\,
  e^{i\bk\cdot\bR},
  \label{eq:Pijkl-fourier}
\end{equation}
where $\tilde{f}_0(\bk) = a^3\,\mathrm{sinc}(k_1 h)\,
\mathrm{sinc}(k_2 h)\,\mathrm{sinc}(k_3 h)/(2\sqrt{\pi})$
is the cube form factor ($h=a/2$) and
$\tilde{\Gamma}_{ijkl}$ is the strain kernel.

The full $9\times 9$ block propagator couples 3~displacement and
6~Voigt strain DOF:
\begin{equation}
  \mathbf{G}_0(\bR,\omega)
  \;=\;
  \begin{pmatrix}
    G_{ij} & C_{i\alpha} \\
    H_{\alpha i} & S_{\alpha\beta}
  \end{pmatrix},
  \label{eq:9x9-structure}
\end{equation}
where $G$~is the $3\times 3$ Green's tensor block,
$S$~is the $6\times 6$ Voigt contraction of~$P_{ijkl}$,
$C$~is the $3\times 6$ displacement--strain coupling
from $\partial\langle G_{ij}\rangle/\partial R_k$,
and $H = C^T$ is its transpose.

This document derives the complete analytical computation of all
blocks for the 26~nearest neighbours on a cubic lattice, with dynamic
corrections through~$\omega^6$.


%======================================================================
\section{The UV Divergence Problem}
\label{sec:uv}
%======================================================================

For the static (Kelvin) Green's tensor, the strain kernel is
\begin{equation}
  \tilde{\Gamma}^\infty_{ijkl}(\kh) \;=\;
  -\frac{1}{2\mu}\bigl[
    \hat{k}_j\hat{k}_l\,\delta_{ik}
  + \hat{k}_i\hat{k}_l\,\delta_{jk}
  - 2\eta\,\hat{k}_i\hat{k}_j\hat{k}_k\hat{k}_l
  \bigr],
  \label{eq:Gamma-inf}
\end{equation}
where $\eta = 1/[2(1-\nu)]$ is the Poisson ratio factor.
This depends only on the \emph{direction}~$\kh$, not on~$|\bk|$.
It is $\order(1)$ as $|\bk|\to\infty$.

The sinc$^2$ form factor provides decay $|\tilde{f}_0|^2 \sim 1/k^2$
per axis, giving $1/k^6$ total.  But the 3D volume element $k^2\,dk$
and the angular measure contribute $k^2$, leaving the radial integral
as $\int dk/k^4 \times k^2 = \int dk/k^2$, which---combined with the
$\order(1)$ kernel---makes the 3D integral~\eqref{eq:Pijkl-fourier}
\textbf{UV-divergent}.  This was confirmed numerically: the integral
does not converge as $K_{\max}\to\infty$.

The resolution is to trade $\kh_j\kh_l = k_j k_l/k^2$ and
$\kh_i\kh_j\kh_k\kh_l = k_ik_jk_kk_l/k^4$ for \emph{real-space
derivatives} of convergent potentials, as described in the next
section.


%======================================================================
\section{Static Propagator via Potential Derivatives}
\label{sec:static}
%======================================================================

\subsection{Reduction to convergent potentials}

The key idea is that $k_j k_l/k^2$ in Fourier space corresponds to
$-\partial^2/\partial R_j\partial R_l$ acting on the $1/k^2$
potential, and similarly $k_ik_jk_kk_l/k^4$ acts on $1/k^4$.
Define:
\begin{align}
  A_{jl}(\bR) &\;=\;
    \int\!\frac{d^3k}{(2\pi)^3}\;
    |\tilde{f}_0|^2\,\hat{k}_j\hat{k}_l\,e^{i\bk\cdot\bR}
  \;=\; -\frac{\partial^2}{\partial R_j\,\partial R_l}
  \biggl[\frac{\Phi_{00}(\bR)}{4\pi}\biggr],
  \label{eq:Ajl}
  \\[4pt]
  B_{ijkl}(\bR) &\;=\;
    \int\!\frac{d^3k}{(2\pi)^3}\;
    |\tilde{f}_0|^2\,\hat{k}_i\hat{k}_j\hat{k}_k\hat{k}_l\,
    e^{i\bk\cdot\bR}
  \;=\; -\frac{\partial^4}{\partial R_i\,\partial R_j\,
    \partial R_k\,\partial R_l}
  \biggl[\frac{\Psi_{00}(\bR)}{8\pi}\biggr],
  \label{eq:Bijkl}
\end{align}
where the \textbf{Newton potential} and \textbf{biharmonic potential}
are tent-weighted integrals:
\begin{align}
  \Phi_{00}(\bR) &\;=\;
  \int M_{00}(w_1-R_1)\,M_{00}(w_2-R_2)\,M_{00}(w_3-R_3)\;
  \frac{1}{|\bw|}\;d^3w,
  \label{eq:Phi00}
  \\[2pt]
  \Psi_{00}(\bR) &\;=\;
  \int M_{00}(w_1-R_1)\,M_{00}(w_2-R_2)\,M_{00}(w_3-R_3)\;
  |\bw|\;d^3w,
  \label{eq:Psi00}
\end{align}
with $M_{00}(d) = \max(0, 1-|d|)$ the tent function
(auto-correlation of two box functions of width~$a$).  Both
integrals converge absolutely.

The static propagator then assembles as:
\begin{equation}
  \boxed{
  P^{(0)}_{ijkl}(\bR) \;=\;
  -\frac{1}{2\mu}\bigl[
    \delta_{ik}\,A_{jl}
  + \delta_{jk}\,A_{il}
  - 2\eta\,B_{ijkl}
  \bigr]
  }
  \label{eq:Pstatic}
\end{equation}
with $\eta = 1/[2(1-\nu)]$.

\subsection{Delta-function collapse}

The derivatives in \eqref{eq:Ajl}--\eqref{eq:Bijkl} act on the tent
functions, not on the kernel.  The second derivative of the tent is a
sum of delta functions:
\begin{equation}
  M_{00}''(d) \;=\; -2\,\delta(d) + \delta(d-1) + \delta(d+1).
  \label{eq:tent-delta}
\end{equation}
Each delta \textbf{collapses one integration dimension} analytically.

\paragraph{Single delta collapse (for $A_{jl}$).}
Differentiating $\Phi_{00}$ twice with respect to~$R_j$ produces
deltas in the $w_j$ integration, leaving a \textbf{2D tent-weighted
master integral}:
\begin{equation}
  I^A_2(p,q,c) \;=\;
  \int_0^1\!\int_0^1
  \frac{u^p\,v^q}{\sqrt{c^2 + u^2 + v^2}}\;du\,dv.
  \label{eq:I2A}
\end{equation}

\paragraph{Double delta collapse (for $B_{ijkl}$).}
When two derivatives act on different tent axes, both dimensions
collapse, leaving a \textbf{1D elementary integral}:
\begin{equation}
  I_1(C) \;=\;
  \int_0^1 (1-t)\,\sqrt{C + t^2}\;dt
  \;=\; \frac{1}{2}\!\left[\sqrt{C+1} + C\,
  \mathrm{arcsinh}\!\left(\frac{1}{\sqrt{C}}\right)\right]
  - \frac{1}{3}\!\left[(C+1)^{3/2} - C^{3/2}\right].
  \label{eq:I1exact}
\end{equation}

\paragraph{Sequential integration.}
All 2D masters are computed analytically by Mathematica via
sequential integration:
\begin{equation}
  \text{inner:}\;
  \int_0^1 \frac{u^p\,v^q}{\sqrt{c^2+u^2+v^2}}\,dv
  \;\;\xrightarrow{\text{Assumptions}}\;\;
  \text{outer:}\;
  \int_0^1 (\cdots)\,du
  \;\;\xrightarrow{\text{Simplify}}\;\;
  \text{exact closed form.}
\end{equation}
Flags: \texttt{Assumptions -> u > 0 \&\& c >= 0},
\texttt{GenerateConditions -> False} at every step.

\subsection{Numerical values}

\begin{table}[ht]
\centering
\caption{Static Newton potential derivatives $A_{jl}$ for unit cube.}
\label{tab:static-A}
\begin{tabular}{l*{4}{r}}
\toprule
Component & Face & Edge & Corner \\
\midrule
$A_{11}$ & $-0.135017$ & $-0.013786$ & $0.0$ (exact) \\
$A_{22}$ & $+0.067509$ & $-0.013786$ & $0.0$ (exact) \\
$A_{33}$ & $+0.067509$ & $+0.027572$ & $0.0$ (exact) \\
$A_{12}$ & $0$ & $-0.045565$ & $-0.016062$ \\
$A_{13}$ & $0$ & $0$ & $-0.016062$ \\
$A_{23}$ & $0$ & $0$ & $-0.016062$ \\
\bottomrule
\end{tabular}
\end{table}

\begin{table}[ht]
\centering
\caption{Static biharmonic potential derivatives $B_{ijkl}$ for unit cube.}
\label{tab:static-B}
\begin{tabular}{l*{3}{r}}
\toprule
Component & Face & Edge & Corner \\
\midrule
$B_{1111}$ & $-0.105182$ & $-0.021889$ & $-0.006255$ \\
$B_{1122}$ & $-0.014918$ & $+0.009657$ & $+0.003128$ \\
$B_{1133}$ & $-0.014918$ & $-0.001554$ & $+0.003128$ \\
$B_{2222}$ & $+0.062267$ & $-0.021889$ & $-0.006255$ \\
$B_{2233}$ & $+0.020159$ & $-0.001554$ & $+0.003128$ \\
$B_{3333}$ & $+0.062267$ & $+0.030680$ & $-0.006255$ \\
$B_{1112}$ & $0$ & $-0.016530$ & $-0.009266$ \\
$B_{1233}$ & $0$ & $-0.012505$ & --- \\
$B_{1123}$ & $0$ & $0$ & $+0.002470$ \\
\bottomrule
\end{tabular}
\end{table}

The Laplacian identity $\sum_k B_{ijkk} = A_{ij}$ is verified to
$10^{-14}$ relative precision for all three neighbour types.


%======================================================================
\section{Three Neighbour Types and $O_h$ Symmetry}
\label{sec:neighbours}
%======================================================================

On a cubic lattice, the 26 nearest neighbours fall into three orbits
under the octahedral group~$O_h$:

\begin{center}
\begin{tabular}{lcccc}
\toprule
Type & Canonical $\bR/a$ & Count & Point group & Independent $A_{jl}$ \\
\midrule
Face   & $(1,0,0)$ & 6  & $C_{4v}$ & $A_{11},\; A_{22}=A_{33}$ \\
Edge   & $(1,1,0)$ & 12 & $C_{2v}$ & $A_{11}=A_{22},\; A_{33},\; A_{12}\neq 0$ \\
Corner & $(1,1,1)$ & 8  & $C_{3v}$ & $A_{11}=A_{22}=A_{33},\; A_{12}=A_{13}=A_{23}$ \\
\bottomrule
\end{tabular}
\end{center}

The independent component counts for each block are:

\begin{center}
\begin{tabular}{l*{4}{c}}
\toprule
Type & $A_{jl}$ & $B_{ijkl}$ & $d^3W/dR^3$ & $dW/dR$ \\
\midrule
Face ($C_{4v}$) & 2 & 4 & 2 & 1 \\
Edge ($C_{2v}$) & 3 & 6 & 3 & 2 \\
Corner ($C_{3v}$) & 2 & 4 & 3 & 1 \\
\bottomrule
\end{tabular}
\end{center}

The delta-function collapse structure differs for each type:
\begin{itemize}[nosep]
  \item \textbf{Face}: one tent shifted by~$a$, two centred.
    Standard 2D masters $I^A_2(p,q,c)$ at $c\in\{0,1,2\}$.
  \item \textbf{Edge}: two tents shifted, one centred.
    Shifted 2D masters
    $I^A_{2,\mathrm{shift}}(p,q,c,s)$ with $\sqrt{c^2+(u+s)^2+v^2}$.
  \item \textbf{Corner}: all three tents shifted.
    Double-shifted 2D masters.  But $C_{3v}$ symmetry reduces the
    count drastically.
\end{itemize}

\paragraph{$O_h$ rotation to all 26 directions.}
Any neighbour direction~$\bR$ is related to the canonical direction
by a signed-permutation matrix $R \in O_h$.  The propagator
transforms as:
\begin{equation}
  P_{ijkl}(\bR) \;=\;
  R_{ia}\,R_{jb}\,R_{kc}\,R_{ld}\,
  P_{abcd}(\bR_{\mathrm{can}}).
  \label{eq:Oh-rotation}
\end{equation}
Only 3~canonical directions need be computed; the remaining
23~follow by symmetry.

\paragraph{Corner diagonal vanishing.}
A notable result is that $A_{11} = A_{22} = A_{33} = 0$
\emph{exactly} for the corner-adjacent case $\bR = (a,a,a)$.
This follows from the identity $B_{1111} = -2B_{1122}$
and the Laplacian relation $B_{1111} + B_{1122} + B_{1133} = A_{11}$.


%======================================================================
\section{Dynamic Corrections: Power Series in $\omega^2$}
\label{sec:dynamic}
%======================================================================

\subsection{Expansion of the dynamic Green's tensor}

The full frequency-dependent Green's tensor in Fourier space is
\begin{equation}
  \tilde{G}_{ij}(\bk,\omega) \;=\;
  \frac{\hat{k}_i\hat{k}_j}%
       {\rho(\omega^2 - c_p^2 k^2)}
  \;+\;
  \frac{\delta_{ij} - \hat{k}_i\hat{k}_j}%
       {\rho(\omega^2 - c_s^2 k^2)}.
  \label{eq:Gij-dynamic}
\end{equation}
Expanding for large~$k$ ($\omega^2/(c^2 k^2) \ll 1$):
\begin{equation}
  \frac{1}{\omega^2 - c^2 k^2}
  \;=\;
  -\frac{1}{c^2 k^2}
  \sum_{n=0}^{\infty}
  \left(\frac{\omega^2}{c^2 k^2}\right)^{\!n}
  \;=\;
  -\sum_{n=0}^{\infty}
  \frac{\omega^{2n}}{c^{2n+2}\,k^{2n+2}}.
  \label{eq:geometric-series}
\end{equation}

The strain kernel has the expansion
\begin{equation}
  \boxed{
  \tilde{\Gamma}_{ijkl}(\bk,\omega)
  \;=\;
  \tilde{\Gamma}^\infty_{ijkl}(\kh)
  \;+\;
  \sum_{n=1}^{\infty}
  \omega^{2n}\,
  \frac{\tilde{\Gamma}^{(n)}_{ijkl}(\kh)}{k^{2n}}
  }
  \label{eq:Gamma-expansion}
\end{equation}
and therefore
\begin{equation}
  P_{ijkl}(\bR,\omega)
  \;=\;
  \sum_{n=0}^{\infty}
  \omega^{2n}\,P^{(n)}_{ijkl}(\bR).
  \label{eq:P-series}
\end{equation}

\subsection{The potential hierarchy}

Each order~$n$ reduces to derivatives of \textbf{higher-order
potentials}:
\begin{equation}
  W_m(\bR) \;=\;
  \int M_{00}(w_1-R_1)\,M_{00}(w_2-R_2)\,M_{00}(w_3-R_3)\;
  |\bw|^{2m-1}\;d^3w,
  \label{eq:Wm-def}
\end{equation}
connected by the Fourier hierarchy
\begin{equation}
  \int\!\frac{d^3k}{(2\pi)^3}\;
  \frac{|\tilde{f}_0|^2}{k^{2m}}\;e^{i\bk\cdot\bR}
  \;=\; c_m\,W_m(\bR).
  \label{eq:FT-hierarchy}
\end{equation}

\begin{center}
\begin{tabular}{ccclc}
\toprule
$m$ & Kernel & Potential & Role & Notation \\
\midrule
$1$ & $1/|\bw|$   & $\Phi$ (Newton)       & Static: $A_{jl}$           & $W_1$ \\
$2$ & $|\bw|$     & $\Psi$ (biharmonic)    & Static: $B_{ijkl}$; Dyn $n$=1 & $W_2$ \\
$3$ & $|\bw|^3$   & $X$ (triharmonic)      & Dyn $n$=1, 2              & $W_3$ \\
$4$ & $|\bw|^5$   & $\Omega$ (pentaharmonic) & Dyn $n$=2, 3            & $W_4$ \\
$5$ & $|\bw|^7$   & $H$ (heptaharmonic)    & Dyn $n$=3, 4              & $W_5$ \\
\bottomrule
\end{tabular}
\end{center}

\subsection{Laplacian chain}

The key relation connecting adjacent potentials is
\begin{equation}
  \nabla^2 |\bw|^{2m+1}
  \;=\; (2m+1)(2m+2)\,|\bw|^{2m-1},
  \label{eq:laplacian-chain}
\end{equation}
giving the \textbf{divisors}:

\begin{center}
\begin{tabular}{cccc}
\toprule
$m$ & Relation & Divisor & Chain \\
\midrule
$1$ & $\nabla^2|\bw| = 2/|\bw|$ & 2 & $\Psi \to \Phi$ \\
$2$ & $\nabla^2|\bw|^3 = 12\,|\bw|$ & 12 & $X \to \Psi$ \\
$3$ & $\nabla^2|\bw|^5 = 30\,|\bw|^3$ & 30 & $\Omega \to X$ \\
$4$ & $\nabla^2|\bw|^7 = 56\,|\bw|^5$ & 56 & $H \to \Omega$ \\
\bottomrule
\end{tabular}
\end{center}

These imply the Laplacian identities for derivatives:
\begin{equation}
  \sum_j \frac{\partial^2 W_{m+1}}{\partial R_j^2}
  \;=\; (2m+1)(2m+2)\,W_m,
  \qquad
  \sum_j B^{(m)}_{ijkk} \;=\; A^{(m)}_{ij}.
  \label{eq:laplacian-identity}
\end{equation}

\subsection{Normalisation constants}

The normalisation constants $c_m$ relating Fourier integrals to
real-space potentials are:

\begin{center}
\begin{tabular}{lll}
\toprule
Constant & Value & Use \\
\midrule
$c_{A,1} = 1/(8\pi)$ & $\approx 0.03979$ & $A^{(1)}$ from $\partial^2\Psi/\partial R^2$ \\
$c_{B,1} = 1/(96\pi)$ & $\approx 0.003316$ & $B^{(1)}$ from $\partial^4 X/\partial R^4$ \\
$c_{A,2} = -1/(96\pi)$ & $\approx -0.003316$ & $A^{(2)}$ from $\partial^2 X/\partial R^2$ \\
$c_{B,2} = -1/(2880\pi)$ & $\approx -0.0001105$ & $B^{(2)}$ from $\partial^4\Omega/\partial R^4$ \\
$c_{A,3} = 1/(2880\pi)$ & $\approx 0.0001105$ & $A^{(3)}$ from $\partial^2\Omega/\partial R^2$ \\
$c_{B,3} = 1/(161280\pi)$ & $\approx 1.969\times 10^{-6}$ & $B^{(3)}$ from $\partial^4 H/\partial R^4$ \\
$c_{A,4} = -1/(161280\pi)$ & $\approx -1.969\times 10^{-6}$ & $A^{(4)}$ from $\partial^2 H/\partial R^2$ \\
\bottomrule
\end{tabular}
\end{center}

Note that $c_{B,n} = c_{A,n}/D_n$ and $c_{A,n+1} = -c_{B,n}$ where
$D_n$ is the divisor at order~$n$.

\subsection{Deviatoric suppression}

At each dynamic order~$n$, the deviatoric (B-tensor) contribution is
suppressed by the factor
\begin{equation}
  \boxed{
  \eta_n \;=\; 1 - \left(\frac{c_s}{c_p}\right)^{2n+2}
  }
  \label{eq:eta-n}
\end{equation}
which approaches unity for large velocity contrast.  The static case
corresponds to $\eta_0 = 1 - (c_s/c_p)^2 = 1/(2(1-\nu))$.

\subsection{Convergence}

The series~\eqref{eq:P-series} converges for
\begin{equation}
  \frac{\omega^2}{c_s^2 k_{\mathrm{eff}}^2} < 1,
  \qquad\text{i.e.}\qquad
  ka \lesssim \pi.
\end{equation}
In the Rayleigh regime ($ka < 0.3$) only $P^{(0)}$ and $P^{(1)}$ are
needed.  In the transition regime ($0.3 < ka < 1$), three or four
terms suffice.

\begin{remark}[Smoothness improves with order]
The Newton potential ($m=1$) has a $1/|\bw|$ singularity---the
hardest integral.  All higher potentials ($m\geq 2$) have smooth
kernels $|\bw|^{2m-1}$.  Successive dynamic corrections are
progressively easier to compute.
\end{remark}


%======================================================================
\section{S Block (Strain--Strain, $6\times 6$)}
\label{sec:S-block}
%======================================================================

The strain--strain propagator at each order has the same tensorial
structure as the static case:
\begin{equation}
  \boxed{
  P^{(n)}_{ijkl}(\bR) \;=\;
  -\frac{1}{2\rho c_s^{2n+2}}\bigl[
    \delta_{ik}\,A^{(n)}_{jl}
  + \delta_{jk}\,A^{(n)}_{il}
  - 2\eta_n\,B^{(n)}_{ijkl}
  \bigr]
  }
  \label{eq:Pn-assembly}
\end{equation}
where $A^{(n)}$ and $B^{(n)}$ are the normalised potential
derivatives at order~$n$:
\begin{center}
\begin{tabular}{cll}
\toprule
Order $n$ & $A^{(n)}_{jl}$ from & $B^{(n)}_{ijkl}$ from \\
\midrule
$0$ (static) & $\partial^2\Phi/(4\pi\partial R^2)$ & $\partial^4\Psi/(8\pi\partial R^4)$ \\
$1$ ($\omega^2$) & $\partial^2\Psi/(8\pi\partial R^2)$ & $\partial^4 X/(96\pi\partial R^4)$ \\
$2$ ($\omega^4$) & $-\partial^2 X/(96\pi\partial R^2)$ & $-\partial^4\Omega/(2880\pi\partial R^4)$ \\
$3$ ($\omega^6$) & $\partial^2\Omega/(2880\pi\partial R^2)$ & $\partial^4 H/(161280\pi\partial R^4)$ \\
\bottomrule
\end{tabular}
\end{center}

The $6\times 6$ Voigt contraction maps $(ij) \to \alpha$ via the
standard ordering: $(0,0)\to 1$, $(1,1)\to 2$, $(2,2)\to 3$,
$(1,2)\to 4$, $(0,2)\to 5$, $(0,1)\to 6$.  Shear columns carry a
factor of~$1/2$ (engineering strain convention).

\subsection{Order 1 ($\omega^2$): $\Psi \to A$, $X \to B$}

\begin{table}[ht]
\centering
\caption{DYN1 raw $A^{(1)}_{jl}$ values ($\times\, c_{A,1} = 1/(8\pi)$).}
\label{tab:dyn1-A}
\begin{tabular}{l*{3}{r}}
\toprule
Component & Face & Edge & Corner \\
\midrule
$A_{11}$ & $0.3043$ & $0.3853$ & $0.3859$ \\
$A_{22}$ & $0.8287$ & $0.3853$ & $0.3859$ \\
$A_{33}$ & $0.8287$ & $0.6464$ & $0.3859$ \\
$A_{12}$ & $0$ & $-0.2679$ & $-0.1612$ \\
\bottomrule
\end{tabular}
\end{table}

\begin{table}[ht]
\centering
\caption{DYN1 raw $B^{(1)}_{ijkl}$ values ($\times\, c_{B,1} = 1/(96\pi)$).}
\label{tab:dyn1-B}
\begin{tabular}{l*{3}{r}}
\toprule
Component & Face & Edge & Corner \\
\midrule
$B_{1111}$ & $1.4133$ & $2.2725$ & $2.5330$ \\
$B_{1122}$ & $1.1194$ & $1.1695$ & $1.0487$ \\
$B_{1133}$ & $1.1194$ & $1.1816$ & $1.0487$ \\
$B_{2222}$ & $6.6139$ & $2.2725$ & $2.5330$ \\
$B_{2233}$ & $2.2113$ & $1.1816$ & $1.0487$ \\
$B_{3333}$ & $6.6139$ & $5.3935$ & $2.5330$ \\
$B_{1112}$ & $0$ & $-1.2918$ & $-0.9196$ \\
$B_{1233}$ & $0$ & $-0.6309$ & --- \\
$B_{1123}$ & $0$ & $0$ & $-0.0954$ \\
\bottomrule
\end{tabular}
\end{table}

\subsection{Order 2 ($\omega^4$): $X \to A$, $\Omega \to B$}

\begin{table}[ht]
\centering
\caption{DYN2 raw $A^{(2)}_{jl}$ values ($\times\, c_{A,2} = -1/(96\pi)$).}
\label{tab:dyn2-A}
\begin{tabular}{l*{3}{r}}
\toprule
Component & Face & Edge & Corner \\
\midrule
$A_{11}$ & $6.1469$ & $6.7280$ & $7.3115$ \\
$A_{22}$ & $3.9310$ & $6.7280$ & $7.3115$ \\
$A_{33}$ & $3.9310$ & $4.9261$ & $7.3115$ \\
$A_{12}$ & $0$ & $1.8087$ & $1.5543$ \\
\bottomrule
\end{tabular}
\end{table}

\begin{table}[ht]
\centering
\caption{DYN2 raw $B^{(2)}_{ijkl}$ values ($\times\, c_{B,2} = -1/(2880\pi)$).}
\label{tab:dyn2-B}
\begin{tabular}{l*{3}{r}}
\toprule
Component & Face & Edge & Corner \\
\midrule
$B_{1111}$ & $121.11$ & $126.69$ & $133.00$ \\
$B_{1122}$ & $31.650$ & $40.522$ & $43.170$ \\
$B_{1133}$ & $31.650$ & $34.625$ & $43.170$ \\
$B_{2222}$ & $64.730$ & $126.69$ & $133.00$ \\
$B_{2233}$ & $21.549$ & $34.625$ & $43.170$ \\
$B_{3333}$ & $64.730$ & $78.533$ & $133.00$ \\
$B_{1112}$ & $0$ & $22.887$ & $20.489$ \\
$B_{1233}$ & $0$ & $8.4875$ & --- \\
$B_{1123}$ & $0$ & $0$ & $5.6510$ \\
\bottomrule
\end{tabular}
\end{table}

\subsection{Order 3 ($\omega^6$): $\Omega \to A$, $H \to B$}

\begin{table}[ht]
\centering
\caption{DYN3 raw $A^{(3)}_{jl}$ values ($\times\, c_{A,3} = 1/(2880\pi)$).}
\label{tab:dyn3-A}
\begin{tabular}{l*{3}{r}}
\toprule
Component & Face & Edge & Corner \\
\midrule
$A_{11}$ & $35.100$ & $51.671$ & $69.712$ \\
$A_{22}$ & $13.542$ & $51.671$ & $69.712$ \\
$A_{33}$ & $13.542$ & $25.500$ & $69.712$ \\
$A_{12}$ & $0$ & $26.145$ & $30.105$ \\
\bottomrule
\end{tabular}
\end{table}

\begin{table}[ht]
\centering
\caption{DYN3 raw $B^{(3)}_{ijkl}$ values ($\times\, c_{B,3} = 1/(161280\pi)$).}
\label{tab:dyn3-B}
\begin{tabular}{l*{3}{r}}
\toprule
Component & Face & Edge & Corner \\
\midrule
$B_{1111}$ & $1401.0$ & $1840.7$ & $2302.4$ \\
$B_{1122}$ & $282.28$ & $651.07$ & $800.71$ \\
$B_{1133}$ & $282.28$ & $401.74$ & $800.71$ \\
$B_{2222}$ & $356.75$ & $1840.7$ & $2302.4$ \\
$B_{2233}$ & $119.31$ & $401.74$ & $800.71$ \\
$B_{3333}$ & $356.75$ & $624.53$ & $2302.4$ \\
$B_{1112}$ & $0$ & $635.51$ & $708.32$ \\
$B_{1233}$ & $0$ & $193.08$ & --- \\
$B_{1123}$ & $0$ & $0$ & $269.26$ \\
\bottomrule
\end{tabular}
\end{table}


%======================================================================
\section{G Block (Displacement--Displacement, $3\times 3$)}
\label{sec:G-block}
%======================================================================

The volume-averaged Green's tensor $\langle G_{ij}\rangle$ follows
the same potential-derivative pattern but with second derivatives
(not fourth):
\begin{equation}
  \boxed{
  \langle G_{ij}\rangle^{(n)} \;=\;
  (-1)^n \frac{\omega^{2n}}{D_n!\;\rho c_s^{2n+2}}
  \biggl[
    \delta_{ij}\,\overline{W}_n
  - \frac{\eta_n}{D_n}\,
    \frac{\partial^2 W_{n+1}}{\partial R_i\,\partial R_j}
  \biggr]
  }
  \label{eq:Gij-assembly}
\end{equation}
where $D_n$ is the Laplacian divisor and
$\overline{W}_n = (1/D_n)\sum_k \partial^2 W_{n+1}/\partial R_k^2$
is the trace (from the Laplacian identity).

\paragraph{Explicit formulae at each order.}
\begin{itemize}[nosep]
\item \textbf{Static} ($n=0$):
  $\langle G_{ij}\rangle^{(0)} =
  \frac{1}{4\pi\mu}\bigl[\delta_{ij}\Phi
  - \frac{1}{4(1-\nu)}\partial^2_{ij}\Psi\bigr]$,
  where $\Phi = \tfrac{1}{2}\sum_k\partial^2_{kk}\Psi$.

\item \textbf{$\omega^2$} ($n=1$):
  $\langle G_{ij}\rangle^{(1)} =
  -\frac{\omega^2}{8\pi\rho c_s^4}\bigl[\delta_{ij}\bar\Psi
  - \frac{\eta_1}{12}\partial^2_{ij}X\bigr]$,
  where $\bar\Psi = \tfrac{1}{12}\sum_k\partial^2_{kk}X$.

\item \textbf{$\omega^4$} ($n=2$):
  $\langle G_{ij}\rangle^{(2)} =
  \frac{\omega^4}{96\pi\rho c_s^6}\bigl[\delta_{ij}\bar X
  - \frac{\eta_2}{30}\partial^2_{ij}\Omega\bigr]$,
  where $\bar X = \tfrac{1}{30}\sum_k\partial^2_{kk}\Omega$.

\item \textbf{$\omega^6$} ($n=3$):
  $\langle G_{ij}\rangle^{(3)} =
  -\frac{\omega^6}{2880\pi\rho c_s^8}\bigl[\delta_{ij}\bar\Omega
  - \frac{\eta_3}{56}\partial^2_{ij}H\bigr]$,
  where $\bar\Omega = \tfrac{1}{56}\sum_k\partial^2_{kk}H$.
\end{itemize}

The $\omega^6$ correction to the G block requires the A-channel
of order~4 (second derivatives of~$H$), but \emph{no} B-tensor
(fourth derivatives of the next potential would be order~5).

\begin{table}[ht]
\centering
\caption{DYN4 raw $A^{(4)}_{jl}$ values for G block $\omega^6$ correction
($\times\, c_{A,4} = -1/(161280\pi)$).}
\label{tab:dyn4-A}
\begin{tabular}{l*{3}{r}}
\toprule
Component & Face & Edge & Corner \\
\midrule
$A_{11}$ & $166.04$ & $336.61$ & $569.49$ \\
$A_{22}$ & $47.160$ & $336.61$ & $569.49$ \\
$A_{33}$ & $47.160$ & $128.50$ & $569.49$ \\
$A_{12}$ & $0$ & $208.61$ & $311.88$ \\
\bottomrule
\end{tabular}
\end{table}


%======================================================================
\section{C/H Blocks (Displacement--Strain Coupling, $3\times 6 + 6\times 3$)}
\label{sec:CH-block}
%======================================================================

The displacement--strain coupling arises from the rank-3 tensor
\begin{equation}
  dG_{ijk}(\bR) \;=\;
  \frac{\partial\langle G_{ij}\rangle}{\partial R_k}
  \label{eq:dG-def}
\end{equation}
which requires \textbf{third derivatives} of the potentials (one
order higher than $A_{jl}$) and \textbf{first derivatives} (one
order lower, obtained via the Laplacian identity).

At each dynamic order:
\begin{equation}
  \boxed{
  dG^{(n)}_{ijk} \;=\;
  (-1)^n \frac{\omega^{2n}}{D_n!\;\rho c_s^{2n+2}}
  \biggl[
    \delta_{ij}\,dW_k
  - \frac{\eta_n}{D_n}\,
    \frac{\partial^3 W_{n+1}}{\partial R_i\,\partial R_j\,\partial R_k}
  \biggr]
  }
  \label{eq:dGijk-assembly}
\end{equation}
where the first derivative $dW_k$ is obtained from the Laplacian:
\begin{equation}
  dW_k \;=\; \frac{1}{D_n}\sum_j
  \frac{\partial^3 W_{n+1}}{\partial R_j^2\,\partial R_k}.
  \label{eq:dW-laplacian}
\end{equation}

The H block ($6\times 3$) is assembled by Voigt contraction:
\begin{equation}
  H_{\alpha i} \;=\; \sum_{jk} V_{\alpha jk}\,dG_{ijk},
  \qquad
  C_{i\alpha} \;=\; \sum_{jk} V_{\alpha jk}\,dG_{jki},
  \qquad
  H = C^T.
  \label{eq:CH-voigt}
\end{equation}

\subsection{Static C/H: $d^3\Psi/dR^3$ and $d\Phi/dR$}

\begin{table}[ht]
\centering
\caption{Biharmonic third derivatives $d^3\Psi/dR_i\,dR_j\,dR_k$ for static C/H.}
\label{tab:d3psi}
\begin{tabular}{l*{3}{r}}
\toprule
Component & Face & Edge & Corner \\
\midrule
$D_{000}$ & $-0.795438$ & $-0.501864$ & $-0.346474$ \\
$D_{011}$ & $-0.528262$ & $+0.048829$ & $-0.021876$ \\
$D_{022}$ & $-0.528262$ & $-0.262771$ & $-0.021876$ \\
$D_{012}$ & --- & --- & $+0.142589$ \\
\bottomrule
\end{tabular}
\end{table}

\begin{table}[ht]
\centering
\caption{Newton first derivatives $d\Phi/dR_k$ (from Laplacian: $d\Phi_k = \tfrac{1}{2}\sum_j D^{\Psi}_{jjk}$).}
\label{tab:dphi}
\begin{tabular}{l*{3}{r}}
\toprule
Component & Face & Edge & Corner \\
\midrule
$d\Phi_0$ & $-0.925981$ & $-0.357903$ & $-0.195113$ \\
$d\Phi_1$ & $0$ & $-0.357903$ & $-0.195113$ \\
$d\Phi_2$ & $0$ & $0$ & $-0.195113$ \\
\bottomrule
\end{tabular}
\end{table}

\subsection{$\omega^2$: $d^3 X/dR^3$ and $d\Psi/dR$}

\begin{table}[ht]
\centering
\caption{Triharmonic third derivatives $d^3 X/dR_i\,dR_j\,dR_k$ ($\omega^2$ C/H correction).}
\label{tab:d3X}
\begin{tabular}{l*{3}{r}}
\toprule
Component & Face & Edge & Corner \\
\midrule
$D_{000}$ & $+5.5907$ & $+4.7979$ & $+4.2544$ \\
$D_{011}$ & $+2.2218$ & $+1.1778$ & $+1.1463$ \\
$D_{022}$ & $+2.2218$ & $+1.8051$ & $+1.1463$ \\
$D_{012}$ & --- & --- & $-0.4109$ \\
\bottomrule
\end{tabular}
\end{table}

\begin{table}[ht]
\centering
\caption{Biharmonic first derivatives $d\Psi/dR_k$
($d\Psi_k = \tfrac{1}{12}\sum_j D^X_{jjk}$).}
\label{tab:dpsi}
\begin{tabular}{l*{3}{r}}
\toprule
Component & Face & Edge & Corner \\
\midrule
$d\Psi_0$ & $+0.836186$ & $+0.648407$ & $+0.545583$ \\
$d\Psi_1$ & $0$ & $+0.648407$ & $+0.545583$ \\
$d\Psi_2$ & $0$ & $0$ & $+0.545583$ \\
\bottomrule
\end{tabular}
\end{table}

\subsection{$\omega^4$: $d^3\Omega/dR^3$ and $dX/dR$}

\begin{table}[ht]
\centering
\caption{Pentaharmonic third derivatives $d^3\Omega/dR_i\,dR_j\,dR_k$ ($\omega^4$ C/H correction).}
\label{tab:d3omega}
\begin{tabular}{l*{3}{r}}
\toprule
Component & Face & Edge & Corner \\
\midrule
$D_{000}$ & $+74.703$ & $+86.935$ & $+97.745$ \\
$D_{011}$ & $+21.532$ & $+34.628$ & $+37.518$ \\
$D_{022}$ & $+21.532$ & $+26.153$ & $+37.518$ \\
$D_{012}$ & --- & --- & $+7.4211$ \\
\bottomrule
\end{tabular}
\end{table}

\begin{table}[ht]
\centering
\caption{Triharmonic first derivatives $dX/dR_k$
($dX_k = \tfrac{1}{30}\sum_j D^\Omega_{jjk}$).}
\label{tab:dX}
\begin{tabular}{l*{3}{r}}
\toprule
Component & Face & Edge & Corner \\
\midrule
$dX_0$ & $+3.9256$ & $+4.9239$ & $+5.7594$ \\
$dX_1$ & $0$ & $+4.9239$ & $+5.7594$ \\
$dX_2$ & $0$ & $0$ & $+5.7594$ \\
\bottomrule
\end{tabular}
\end{table}

\subsection{$\omega^6$: $d^3 H/dR^3$ and $d\Omega/dR$}

\begin{table}[ht]
\centering
\caption{Heptaharmonic third derivatives $d^3 H/dR_i\,dR_j\,dR_k$ ($\omega^6$ C/H correction).}
\label{tab:d3H}
\begin{tabular}{l*{3}{r}}
\toprule
Component & Face & Edge & Corner \\
\midrule
$D_{000}$ & $+520.37$ & $+818.18$ & $+1154.5$ \\
$D_{011}$ & $+119.39$ & $+401.75$ & $+531.43$ \\
$D_{022}$ & $+119.39$ & $+208.56$ & $+531.43$ \\
$D_{012}$ & --- & --- & $+219.48$ \\
\bottomrule
\end{tabular}
\end{table}

\begin{table}[ht]
\centering
\caption{Pentaharmonic first derivatives $d\Omega/dR_k$
($d\Omega_k = \tfrac{1}{56}\sum_j D^H_{jjk}$).}
\label{tab:domega}
\begin{tabular}{l*{3}{r}}
\toprule
Component & Face & Edge & Corner \\
\midrule
$d\Omega_0$ & $+13.556$ & $+25.509$ & $+39.595$ \\
$d\Omega_1$ & $0$ & $+25.509$ & $+39.595$ \\
$d\Omega_2$ & $0$ & $0$ & $+39.595$ \\
\bottomrule
\end{tabular}
\end{table}


%======================================================================
\section{Validation and Summary}
\label{sec:validation}
%======================================================================

\subsection{Laplacian identities}

The fundamental check is the Laplacian identity
\begin{equation}
  \sum_k B^{(n)}_{ijkk} \;=\; A^{(n)}_{ij}
  \label{eq:laplacian-check}
\end{equation}
at every dynamic order for all three neighbour types.  This is
verified to $10^{-9}$--$10^{-12}$ relative precision.

For the C/H blocks, the analogous identity is
\begin{equation}
  \sum_j \frac{\partial^3 W_{m+1}}{\partial R_j^2\,\partial R_k}
  \;=\; D_m \cdot \frac{\partial W_m}{\partial R_k}
  \label{eq:laplacian-ch}
\end{equation}
with divisors $D_m \in \{2, 12, 30, 56\}$.  This is verified to
18+ digits for the static case and 12+ digits for dynamic orders.

\subsection{Point-group symmetry}

All tensor components respect the appropriate point-group symmetry:
\begin{itemize}[nosep]
  \item \textbf{Face} ($C_{4v}$): $A_{22}=A_{33}$, $A_{12}=A_{13}=A_{23}=0$,
    $B_{2222}=B_{3333}$, $B_{1122}=B_{1133}$, all mixed $B=0$.
  \item \textbf{Edge} ($C_{2v}$): $A_{11}=A_{22}$, $A_{13}=A_{23}=0$,
    $B_{1111}=B_{2222}$, $B_{1133}=B_{2233}$, $B_{1112}=B_{1222}$,
    $B_{1123}=B_{2213}=0$.
  \item \textbf{Corner} ($C_{3v}$): full $S_3$ permutation symmetry.
    $A_{11}=A_{22}=A_{33}=0$ (exact), all off-diagonal $A$ equal,
    $B_{1111}=-2B_{1122}$.
\end{itemize}

\subsection{$H = C^T$ preservation}

The identity $H_{\alpha i} = C_{i\alpha}$ (transpose symmetry of the
coupling blocks) is preserved at all dynamic orders.  This follows
from the symmetry of $dG_{ijk}$ in its first two indices (inherited
from $G_{ij} = G_{ji}$).

\subsection{Inversion symmetry}

Under $\bR \to -\bR$:
\begin{itemize}[nosep]
  \item $G_{ij}(-\bR) = G_{ij}(\bR)$ --- even (second derivatives of
    even potentials).
  \item $C_{i\alpha}(-\bR) = -C_{i\alpha}(\bR)$ --- odd (first/third
    derivatives pick up a sign).
  \item $S_{\alpha\beta}(-\bR) = S_{\alpha\beta}(\bR)$ --- even
    (fourth derivatives of even potentials).
\end{itemize}

\subsection{Convergence}

Each successive dynamic order is smaller than the previous one
for $ka < 1$:
\begin{equation}
  \frac{|P^{(n+1)}|}{|P^{(n)}|}
  \;\sim\;
  \left(\frac{ka}{\pi}\right)^{\!2}
  \;\ll 1
  \quad\text{for}\quad ka < 1.
\end{equation}

\subsection{Dynamic coverage summary}

\begin{center}
\begin{tabular}{lccccc}
\toprule
Block & Size & $\omega^0$ & $\omega^2$ & $\omega^4$ & $\omega^6$ \\
\midrule
S (strain--strain)   & $6\times 6$ & \checkmark & \checkmark & \checkmark & \checkmark \\
G (disp.--disp.)     & $3\times 3$ & \checkmark & \checkmark & \checkmark & \checkmark \\
C (disp.--strain)    & $3\times 6$ & \checkmark & \checkmark & \checkmark & \checkmark \\
H (strain--disp.)    & $6\times 3$ & \checkmark & \checkmark & \checkmark & \checkmark \\
\bottomrule
\end{tabular}
\end{center}

All four blocks of the $9\times 9$ propagator are complete through
$\omega^6$, for all 26 nearest neighbours.  The total number of
precomputed analytical constants is approximately~200.

\subsection{Potential hierarchy}

The potentials used at each block and order are:

\begin{center}
\begin{tabular}{l*{5}{c}}
\toprule
& $\Phi$ & $\Psi$ & $X$ & $\Omega$ & $H$ \\
& ($1/\rho$) & ($\rho$) & ($\rho^3$) & ($\rho^5$) & ($\rho^7$) \\
\midrule
S: $\partial^2/\partial R^2$ (A) & $\omega^0$ & $\omega^2$ & $\omega^4$ & $\omega^6$ & --- \\
S: $\partial^4/\partial R^4$ (B) & --- & $\omega^0$ & $\omega^2$ & $\omega^4$ & $\omega^6$ \\
G: $\partial^2/\partial R^2$     & --- & $\omega^0$ & $\omega^2$ & $\omega^4$ & $\omega^6$ \\
C/H: $\partial^3/\partial R^3$   & --- & $\omega^0$ & $\omega^2$ & $\omega^4$ & $\omega^6$ \\
C/H: $\partial/\partial R$       & $\omega^0$ & $\omega^2$ & $\omega^4$ & $\omega^6$ & --- \\
\bottomrule
\end{tabular}
\end{center}

No numerical quadrature is needed at any stage.  Python consumes the
precomputed constants for the Foldy--Lax solver.

\end{document}
